\documentclass[12pt, a4paper]{article}
\usepackage[T2A]{fontenc}
\usepackage[utf8]{inputenc}
\usepackage[russian]{babel}
\usepackage{amsmath, amssymb} % Для математических символов
\usepackage{enumitem} % Для гибкой настройки списков
\usepackage{geometry} % Для настройки полей
\geometry{top=2cm, bottom=2cm, left=2.5cm, right=2cm}

\begin{document}

\section*{Георгий Игоревич Вольфсон <<Делимость с человеческим лицом>>}

\subsection*{Беседа третья. Простые и составные числа. Делители числа.}

\subsubsection*{Упражнения}

\begin{enumerate}[label=\arabic*., wide=0pt, leftmargin=*]

\item \textbf{Простыми или составными являются числа: а) 45, б) 92, в) 109, г) 121, д) 343, е) 463, ж) 3599?}

\textbf{Решение:}
\begin{enumerate}[label=\alph*), wide=0pt, leftmargin=*]
\item $45 = 5 \times 9$, следовательно, составное.
\item $92 = 2 \times 46$ (или $92 = 4 \times 23$), следовательно, составное.
\item $109$ — простое (проверено на делимость на 2, 3, 5, 7, 11: ни одно из них не делит 109).
\item $121 = 11^2$, следовательно, составное.
\item $343 = 7^3$, следовательно, составное.
\item $463$ — простое (проверено на делимость на 2, 3, 5, 7, 11, 13, 17, 19; $19^2 = 361$, $23^2 = 529 > 463$, дальнейшая проверка не требуется).
\item $3599 = 3600 - 1 = 60^2 - 1 = (60-1)(60+1) = 59 \times 61$, следовательно, составное.
\end{enumerate}

\item \textbf{При каких натуральных $n$ число $3n$ является простым?}

\textbf{Решение:} 
При $n = 1$ получаем $3 \times 1 = 3$ — простое число. 
При всех $n > 1$ число $3n$ имеет делители $1$, $3$ и $3n$, а значит является составным. 
Таким образом, ответ: $n = 1$.

\item \textbf{При каких натуральных $n$ число $n^2 - 1$ является простым?}

\textbf{Решение:} 
Разложим на множители: $n^2 - 1 = (n - 1)(n + 1)$. 
Для того чтобы произведение двух натуральных множителей было простым, необходимо, чтобы один из множителей равнялся $1$. 
\begin{itemize}
\item Если $n - 1 = 1$, то $n = 2$, и тогда $n + 1 = 3$, $n^2 - 1 = 3$ — простое.
\item Если $n + 1 = 1$, то $n = 0$, что не является натуральным числом.
\end{itemize}
Проверка: при $n = 2$ число $n^2 - 1 = 3$ действительно простое. 
Следовательно, ответ: $n = 2$.

\item \textbf{Докажите, что $x^2 + 8x + 7$ не является простым числом ни при каком натуральном $x$.}

\textbf{Решение:} 
Разложим многочлен на множители:
\[
x^2 + 8x + 7 = (x+1)(x+7).
\]
Для того чтобы произведение двух натуральных чисел было простым, необходимо, чтобы один из множителей равнялся $1$.
\begin{itemize}
\item Если $x+1 = 1$, то $x = 0$, что не является натуральным числом.
\item Если $x+7 = 1$, то $x = -6$, что также не является натуральным.
\end{itemize}
Следовательно, ни при каком натуральном $x$ произведение $(x+1)(x+7)$ не может быть простым числом. 
При $x \ge 1$ оба множителя больше $1$, и число имеет нетривиальные делители.

\item \textbf{При каких натуральных $x$ число $x^2 + 4x$ является простым?}

\textbf{Решение:} 
Разложим на множители:
\[
x^2 + 4x = x(x+4).
\]
Рассмотрим возможности:
\begin{itemize}
\item Если $x = 1$, то $x+4 = 5$, и произведение $1 \times 5 = 5$ — простое число.
\item Если $x+4 = 1$, то $x = -3$, что не является натуральным.
\end{itemize}
При $x > 1$ оба множителя $x$ и $x+4$ больше $1$, поэтому число составное.
Ответ: $x = 1$.

\item \textbf{Запишите каноническое разложение числа: а) 250, б) 1200, в) 2424.}

\textbf{Решение:}
\[
\begin{aligned}
&\text{а)}\ 250 = 2 \times 125 = 2^1 \times 5^3,\\
&\text{б)}\ 1200 = 12 \times 100 = (2^2 \times 3) \times (2^2 \times 5^2) = 2^4 \times 3^1 \times 5^2,\\
&\text{в)}\ 2424 = 8 \times 303 = 2^3 \times (3 \times 101) = 2^3 \times 3^1 \times 101^1.
\end{aligned}
\]

\item \textbf{Найдите НОД и НОК следующих чисел: а) (90; 64); б) (250; 512); в) (900; 3000).}

\textbf{Решение:}
\[
\begin{aligned}
&\text{а)}\ 90 = 2 \times 3^2 \times 5,\quad 64 = 2^6.\\
&\text{НОД} = 2,\quad \text{НОК} = 2^6 \times 3^2 \times 5 = 64 \times 9 \times 5 = 2880.\\[4pt]
&\text{б)}\ 250 = 2 \times 5^3,\quad 512 = 2^9.\\
&\text{НОД} = 2,\quad \text{НОК} = 2^9 \times 5^3 = 512 \times 125 = 64000.\\[4pt]
&\text{в)}\ 900 = 2^2 \times 3^2 \times 5^2,\quad 3000 = 2^3 \times 3 \times 5^3.\\
&\text{НОД} = 2^2 \times 3 \times 5^2 = 4 \times 3 \times 25 = 300,\\
&\text{НОК} = 2^3 \times 3^2 \times 5^3 = 8 \times 9 \times 125 = 9000.
\end{aligned}
\]

\item \textbf{Найдите количество делителей числа: а) 7; б) 15; в) 24; г) 2000.}

\textbf{Решение:} 
Если каноническое разложение числа $n = p_1^{\alpha_1} p_2^{\alpha_2} \dots p_k^{\alpha_k}$, то количество делителей равно $(\alpha_1+1)(\alpha_2+1)\dots(\alpha_k+1)$.
\[
\begin{aligned}
&\text{а)}\ 7 = 7^1 \quad\Rightarrow\quad 1+1 = 2\ \text{делителя (1 и 7)}.\\
&\text{б)}\ 15 = 3^1 \times 5^1 \quad\Rightarrow\quad (1+1)(1+1) = 4\ \text{делителя}.\\
&\text{в)}\ 24 = 2^3 \times 3^1 \quad\Rightarrow\quad (3+1)(1+1) = 8\ \text{делителей}.\\
&\text{г)}\ 2000 = 2^4 \times 5^3 \quad\Rightarrow\quad (4+1)(3+1) = 5 \times 4 = 20\ \text{делителей}.
\end{aligned}
\]

\item \textbf{Найдите сумму делителей числа: а) 13; б) 28; в) 64; г) 300.}

\textbf{Решение:} 
Сумма делителей вычисляется по формуле:
\[
\sigma(n) = \frac{p_1^{\alpha_1+1}-1}{p_1-1} \cdot \frac{p_2^{\alpha_2+1}-1}{p_2-1} \cdots
\]
\[
\begin{aligned}
&\text{а)}\ 13 = 13^1 \quad\Rightarrow\quad \frac{13^2-1}{13-1} = \frac{169-1}{12} = \frac{168}{12} = 14\ (1+13).\\
&\text{б)}\ 28 = 2^2 \times 7^1 \quad\Rightarrow\quad \frac{2^3-1}{2-1} \cdot \frac{7^2-1}{7-1} = \frac{8-1}{1} \cdot \frac{49-1}{6} = 7 \cdot \frac{48}{6} = 7 \cdot 8 = 56.\\
&\text{в)}\ 64 = 2^6 \quad\Rightarrow\quad \frac{2^7-1}{2-1} = \frac{128-1}{1} = 127.\\
&\text{г)}\ 300 = 2^2 \times 3^1 \times 5^2 \quad\Rightarrow\quad 
\frac{2^3-1}{1} \cdot \frac{3^2-1}{2} \cdot \frac{5^3-1}{4} = \\ 
&\qquad\qquad\qquad\qquad\qquad\; = 7 \cdot \frac{9-1}{2} \cdot \frac{125-1}{4} = 
7 \cdot \frac{8}{2} \cdot \frac{124}{4} = 7 \cdot 4 \cdot 31 = 868.
\end{aligned}
\]

\item \textbf{Найдите НОД и НОК следующих чисел: а) 24, 30, 42; б) 100, 300, 900; в) 16, 18, 20, 22.}

\textbf{Решение:}

\begin{enumerate}[label=\alph*), wide=0pt, leftmargin=*]
\item 
\(24 = 2^3 \times 3,\quad 30 = 2 \times 3 \times 5,\quad 42 = 2 \times 3 \times 7.\)

НОД (берём минимальные степени общих простых): \(2^1 \times 3^1 = 6.\)

НОК (берём максимальные степени всех простых): \(2^3 \times 3^1 \times 5^1 \times 7^1 = 8 \times 3 \times 5 \times 7 = 840.\)

\item 
\(100 = 2^2 \times 5^2,\quad 300 = 2^2 \times 3 \times 5^2,\quad 900 = 2^2 \times 3^2 \times 5^2.\)

НОД: \(2^2 \times 5^2 = 4 \times 25 = 100.\)

НОК: \(2^2 \times 3^2 \times 5^2 = 4 \times 9 \times 25 = 900.\)

\item 
\(16 = 2^4,\quad 18 = 2 \times 3^2,\quad 20 = 2^2 \times 5,\quad 22 = 2 \times 11.\)

НОД: общий делитель для всех чисел — только \(2^1 = 2.\)

НОК: \(2^4 \times 3^2 \times 5^1 \times 11^1 = 16 \times 9 \times 5 \times 11 = 144 \times 55 = 7920.\)
\end{enumerate}

\item \textbf{Может ли у круглого числа быть ровно а) 7, б) 8 натуральных делителей?}

\textbf{Решение:} 
Круглое число — это число, оканчивающееся на 0, то есть имеющее в каноническом разложении множители \(2\) и \(5\) как минимум в первой степени: \(n = 2^{\alpha} \cdot 5^{\beta} \cdot m\), где \(\alpha \ge 1\), \(\beta \ge 1\), а \(m\) не делится на 2 и 5.

Количество делителей: \(\tau(n) = (\alpha+1)(\beta+1) \cdot \tau(m)\), причём \(\tau(m) \ge 1\).

\begin{enumerate}[label=\alph*), wide=0pt, leftmargin=*]
\item Для 7 делителей: 7 — простое число. Значит, \(\tau(n) = 7\) возможно только если \((\alpha+1)(\beta+1) = 7\) и \(\tau(m)=1\). Но 7 простое, поэтому один из сомножителей равен 7, другой — 1. Однако \(\alpha+1 \ge 2\) и \(\beta+1 \ge 2\), их произведение не может быть 7. Следовательно, 7 делителей быть не может.

\item Для 8 делителей: например, \(\alpha=1,\ \beta=1,\ m=1\) даёт \(\tau = (1+1)(1+1) = 4\) — не подходит. Подберём: \(\alpha=3,\ \beta=1,\ m=1\) даёт \((3+1)(1+1)=8\). То есть число \(2^3 \times 5^1 = 8 \times 5 = 40\) имеет ровно 8 делителей. Проверка: 40 — круглое число. Ответ: 8 делителей может быть (например, 40).
\end{enumerate}
Таким образом: а) нет; б) да.

\item \textbf{Гоша вычислил сумму всех делителей некоторой степени двойки и получил 1 234 567 890. Докажите, что он ошибся.}

\textbf{Решение:}
Сумма всех делителей числа \(2^k\) (где \(k \ge 0\)) равна:
\[
\sigma(2^k) = 1 + 2 + 4 + \dots + 2^k = 2^{k+1} - 1.
\]
Это число всегда нечётное, так как \(2^{k+1}\) чётно, а \(2^{k+1} - 1\) нечётно.
Однако \(1\,234\,567\,890\) — число чётное (оканчивается на 0). Следовательно, оно не может быть равно \(2^{k+1}-1\) ни при каком натуральном \(k\). Значит, Гоша ошибся.

\item \textbf{Докажите, что произведение любых пяти последовательных натуральных чисел делится на 30.}

\textbf{Решение:}
Рассмотрим произведение \(n(n+1)(n+2)(n+3)(n+4)\).

\begin{itemize}
\item Среди любых пяти последовательных чисел есть хотя бы одно число, кратное 5.
\item Среди любых пяти последовательных чисел есть хотя бы одно число, кратное 3.
\item Среди любых пяти последовательных чисел есть хотя бы два чётных числа, одно из них кратно 4, поэтому в произведении есть множитель \(2^2\).
\end{itemize}
Таким образом, в каноническом разложении произведения присутствуют множители \(2^2\), \(3\) и \(5\). Следовательно, произведение делится на \(2^2 \times 3 \times 5 = 60\), а значит, и на \(30\).

\end{enumerate}

\end{document}